\input{preamble.tex}

\title{Лабораторная работа № 6. \\ Статическая маршрутизация VLAN}
\author{Сущенко Алина Николаевна \\ НПИбд-01-23}

\institute{Российский университет дружбы народов имени Патриса Лумумбы}
\date{2026}

\begin{document}

\frame{\titlepage}

\begin{frame}
\frametitle{Цель работы}
\begin{itemize}
    \item Настроить статическую маршрутизацию VLAN в сети.
\end{itemize}
\end{frame}

\begin{frame}
\frametitle{Топология сети}
    \centering
    \includegraphics[width=0.9\textwidth]{../images/01.png}
    \captionof{figure}{Обновлённая топология сети.}
\end{frame}

\begin{frame}[containsverbatim]
\frametitle{Первоначальная настройка маршрутизатора}
    \centering
    \includegraphics[width=0.8\textwidth]{../images/02.png}
    \captionof{figure}{Первоначальная настройка маршрутизатора.}
\end{frame}

\begin{frame}[containsverbatim]
\frametitle{Настройка виртуальных интерфейсов на маршрутизаторе}
    \centering
    \includegraphics[width=0.9\textwidth]{../images/03.png}
    \captionof{figure}{Настройка виртуальных интерфейсов, соответствующих номерам VLAN.}
\end{frame}

\begin{frame}
\frametitle{Настройка trunk-портов}
    \centering
    \includegraphics[width=0.9\textwidth]{../images/05.png}
    \captionof{figure}{Настройка trunk-портов.}
\end{frame}

\begin{frame}
\frametitle{Проверка настроек маршрутизатора}
    \centering
    \includegraphics[width=0.9\textwidth]{../images/04.png}
    \captionof{figure}{Проверка заданных настроек.}
\end{frame}

\begin{frame}
\frametitle{Результаты проверки}
\begin{itemize}
    \item Физический интерфейс FastEthernet0/0: up/up
    \item Все суб-интерфейсы (VLAN 2,3,101,102,103,104) в состоянии up/up с правильными IP-адресами
    \item Команда \texttt{show vlan} на маршрутизаторе показывает только стандартные VLAN, что является нормой
    \item Команда \texttt{show interfaces trunk} не показывает информацию на маршрутизаторе
\end{itemize}
\end{frame}

\begin{frame}
\frametitle{Проверка связи ПК dk-ansusenko-01 (VLAN 101)}
    \centering
    \includegraphics[width=0.9\textwidth]{../images/06.png}
    \captionof{figure}{Проверка связи с ПК dk-ansusenko-01.}
\end{frame}

\begin{frame}
\frametitle{Результаты проверки VLAN 101}
\begin{itemize}
    \item \textbf{Пинг шлюза (10.128.3.1):} успешно (TTL=255) - связь в пределах VLAN работает
    \item \textbf{Пинг ПК из VLAN 102 (10.128.4.2):} успешно (TTL=127) - межVLAN маршрутизация работает
    \item \textbf{Пинг сервера (10.128.0.2):} успешно (TTL=127) - доступ к серверной ферме работает
\end{itemize}
\end{frame}

\begin{frame}
\frametitle{Проверка связи ПК из VLAN 102}
    \centering
    \includegraphics[width=0.9\textwidth]{../images/07.png}
    \captionof{figure}{Проверка связи ПК из VLAN 102.}
\end{frame}

\begin{frame}
\frametitle{Проверка связи между серверами}
    \centering
    \includegraphics[width=0.9\textwidth]{../images/08.png}
    \captionof{figure}{Проверка связи между серверами.}
\end{frame}

\begin{frame}
\frametitle{Режим симуляции в Packet Tracer}
    \centering
    \includegraphics[width=0.9\textwidth]{../images/09.png}
    \captionof{figure}{Просмотр пакетов после завершения симуляции.}
\end{frame}

\begin{frame}
\frametitle{Таблица распределения IP-адресов}
\begin{table}[h]
    \centering
    \tiny
    \begin{tabular}{|p{2.5cm}|c|c|c|}
        \hline
        \textbf{Устройство} & \textbf{VLAN} & \textbf{IP-адрес} & \textbf{Назначение} \\
        \hline
        msk-ansusenko-gw-1 & - & 10.128.0.1 & Шлюз серверной фермы \\
        \hline
        msk-ansusenko-gw-1 & - & 10.128.1.1 & Шлюз управления \\
        \hline
        msk-ansusenko-gw-1 & - & 10.128.3.1 & Шлюз ДК \\
        \hline
        msk-ansusenko-gw-1 & - & 10.128.4.1 & Шлюз кафедр \\
        \hline
        msk-ansusenko-gw-1 & - & 10.128.5.1 & Шлюз администрации \\
        \hline
        msk-ansusenko-gw-1 & - & 10.128.6.1 & Шлюз других пользователей \\
        \hline
        ansusenko-web & 3 & 10.128.0.2 & Web-сервер \\
        \hline
        ansusenko-file & 3 & 10.128.0.3 & Файловый сервер \\
        \hline
        ansusenko-mail & 3 & 10.128.0.4 & Почтовый сервер \\
        \hline
        dk-ansusenko-01 & 101 & 10.128.3.2 & ПК пользователя (ДК) \\
        \hline
        dep-ansusenko-1 & 102 & 10.128.4.2 & ПК пользователя (кафедра) \\
        \hline
        adm-ansusenko-1 & 103 & 10.128.5.2 & ПК пользователя (администрация) \\
        \hline
        other-ansusenko-1 & 104 & 10.128.6.2 & ПК пользователя (другие) \\
        \hline
    \end{tabular}
\end{table}
\end{frame}

\begin{frame}

\frametitle{Контрольные вопросы}

\begin{block}{1. Охарактеризуйте стандарт IEEE 802.1Q}
\begin{itemize}
    \item Открытый стандарт для создания виртуальных локальных сетей (VLAN)
    \item Определяет механизм тегирования (маркировки) Ethernet-кадров
    \item Позволяет передавать трафик нескольких VLAN по одному физическому каналу (trunk)
    \item Обеспечивает изоляцию трафика на канальном уровне модели OSI
    \item Поддерживается оборудованием разных производителей
\end{itemize}
\end{block}

\begin{block}{2. Опишите формат кадра IEEE 802.1Q}
\begin{itemize}
    \item Вставляет 4-байтовый тег между MAC-адресом источника и полем Type/Length
    \item \textbf{TPID (Tag Protocol ID - 2 байта):} значение 0x8100, указывает на наличие тега VLAN
    \item \textbf{TCI (Tag Control Information - 2 байта):}
    \begin{itemize}
        \item Priority (3 бита) - приоритет кадра (CoS)
        \item CFI (1 бит) - индикатор формата кадра
        \item VID (12 бит) - идентификатор VLAN (1-4094)
    \end{itemize}
    \item После тега следует оригинальное поле Type/Length
    \item CRC пересчитывается с учетом добавленного тега
\end{itemize}
\end{block}

\end{frame}

\begin{frame}
\frametitle{Выводы}
\begin{itemize}
    \item В результате выполнения лабораторной работы мы построили статическую маршрутизацию VLAN.
\end{itemize}
\end{frame}

\end{document}
